% Extracted from p03-gqa-small
\subsection{Grouped-Query Attention at Small Scale}
\label{ch-gqa}

\paragraph{Background}

\paragraph{Attention variants.} With model width $d$, $h_q$ query
heads and head dimension $d_h = d/h_q$: multi-head attention (MHA)
uses $h_{kv} = h_q$ key/value heads; multi-query attention uses
$h_{kv} = 1$; grouped-query attention \citep{ainslie2023gqa} uses
$1 < h_{kv} < h_q$, with each KV head shared by $h_q/h_{kv}$ query
heads. GQA was shown to match MHA quality closely while
substantially beating single-head sharing on quality.

\paragraph{Parameter arithmetic}
\label{sec:params}

Per attention layer, with biases ignored:
\begin{equation}
  P_{\mathrm{MHA}} = \underbrace{d^2}_{Q} + \underbrace{d^2}_{K}
  + \underbrace{d^2}_{V} + \underbrace{d^2}_{O} = 4d^2.
\end{equation}
For GQA with $h_{kv}$ KV heads, the K and V projections map
$d \to h_{kv} d_h$:
\begin{equation}
  P_{\mathrm{GQA}}(h_{kv}) = 2d^2 + 2\,d\,h_{kv}d_h
  = 2d^2\left(1 + \frac{h_{kv}}{h_q}\right).
\end{equation}
With $d = 512$, $h_q = 8$, $d_h = 64$: MHA costs
$4 \times 262{,}144 = 1{,}048{,}576$ parameters per layer. GQA with
$h_{kv} = 2$ costs $524{,}288 + 131{,}072 = 655{,}360$ per layer, a
saving of $393{,}216$ parameters per layer ($37.5\%$ of the
attention sub-layer). Across the 8 unique attention layers of our
reference model (6 prelude $+$ 2 coda), that is $\approx 3.15$M
parameters, or about $7.2\%$ of the 43,974,722 total. A welcome
discount, but not the reason to adopt GQA.

\paragraph{KV-cache arithmetic}
\label{sec:cache}

The KV cache stores, per token per layer, one key and one value
vector per KV head. In fp16:
\begin{equation}
  C_{\mathrm{tok,layer}}(h_{kv}) = 2 \cdot h_{kv} \cdot d_h
  \cdot 2\ \mathrm{bytes}.
\end{equation}
For our dimensions: MHA ($h_{kv} = 8$) needs $2{,}048$ bytes per
token per layer; GQA with $h_{kv} = 2$ needs $512$ bytes---a factor
$h_q/h_{kv} = 4$ reduction, independent of model size. Total cache:
\begin{equation}
  C_{\mathrm{total}} = U \cdot C_{\mathrm{tok,layer}} \cdot T
  \cdot B,
\end{equation}
where $U$ is the number of unrolled layers holding cache, $T$ the
context length and $B$ the batch size. Note that in a looped
architecture, each loop iteration produces its own KV entries, so
$U$ counts unrolled depth (prelude $+$ loops $\times$ loop block
$+$ coda), not unique layers.

\paragraph{Worked serving scenarios}
\label{sec:scenarios}

Take $U = 12$ unrolled layers as an illustrative depth and fp16
throughout. Model weights occupy $\approx 88$MB. All sizes below are
decimal (1\,MB $= 10^6$ bytes).

\begin{center}
{\small\setlength{\tabcolsep}{4pt}
\begin{tabular}{lrrrr}
\toprule
Scenario & $B$ & $T$ & MHA cache & GQA-2 cache \\
\midrule
Single chat & 1 & 512 & 12.6\,MB & 3.1\,MB \\
Batched serving & 8 & 2{,}048 & 403\,MB & 101\,MB \\
Heavy batch & 32 & 4{,}096 & 3.22\,GB & 0.81\,GB \\
\bottomrule
\end{tabular}}
\end{center}

The pattern is stark. For a single short conversation, the MHA
cache (12.6\,MB) is a footnote next to 88\,MB of weights---GQA buys
almost nothing. At batch 8 and 2k context, the MHA cache
(403\,MB) is over $4\times$ the weights, and GQA reclaims some
300\,MB. At heavy batching it is the difference between fitting on
one modest GPU and not fitting at all.

There is a second, throughput-shaped win: autoregressive decoding
is memory-bandwidth bound, and each step streams the whole cache.
GQA cuts bytes-per-step by $h_q/h_{kv}$, which translates directly
into higher token throughput in cache-bound regimes---again, only
where the cache is large enough to be the bottleneck.

\paragraph{Why training barely notices}
\label{sec:training}

During training there is no KV cache; the K and V \emph{activations}
are recomputed per forward pass. GQA shrinks them by $h_{kv}/h_q$,
but K/V activations are a small slice of total activation memory,
which is dominated by $O(BTd)$ terms (Q, O, FFN intermediates,
residuals) that GQA does not touch. Expect low-single-digit percent
training memory savings and no meaningful throughput change: GQA is
a serving optimisation, full stop.

\paragraph{Discussion: when GQA matters at small scale}

The decision rule falls out of Section~\ref{sec:scenarios}:
compute $C_{\mathrm{total}}$ for MHA in your serving regime; if it
is comparable to or larger than the weight footprint, GQA (or even
fewer KV heads) is worth it; if it is an order of magnitude
smaller, it is cosmetic. Concretely for 44M--500M models:
\begin{itemize}
  \item \textbf{On-device / single-user chat:} negligible. Skip it
    unless the parameter saving is wanted for its own sake.
  \item \textbf{Batched inference servers, long context:}
    decisive. The cache dominates and GQA's $4\times$ cut (at
    $h_{kv} = 2$) directly buys batch size and context length.
  \item \textbf{Speculative or parallel decoding:} any technique
    that multiplies effective batch size pushes you into the
    cache-dominated regime sooner.
\end{itemize}
The quality cost, per \citet{ainslie2023gqa}, is small at these
grouping ratios, so the trade is nearly free where it matters and
unnecessary where it does not.

\paragraph{Limitations}

All figures assume fp16 and the stated head geometry; int8/int4
quantised caches scale all numbers down uniformly without changing
the ratios. We analyse memory, not the engineering cost of
implementing grouped attention kernels, which is real but
one-off. Quality parity is taken from \citet{ainslie2023gqa}
rather than re-measured here.

\paragraph{Conclusion}

GQA's value proposition does not fade at small scale---it
\emph{relocates}. The parameter saving ($\sim$7\% at 44M) is a
bonus; the KV-cache factor of $h_q/h_{kv}$ is the main event, and
it matters exactly when serving memory is cache-bound: large
batches, long contexts. For short-context single-user inference it
is negligible, and for training it is nearly invisible. Size the
cache for your regime, then decide.
